\documentstyle[% 


righttag% 


,11pt% 


,amssymb% 


,russian%               remove in Eng. 


,izvru%                 Set izven% in Eng. 


%,izvru-ks%            for brief communications 


]{amsart} 


 


% 


\newcommand{\tts}[1]{} 


 


%\hoffset=-15mm%                  For Eng. 


\setcounter{page}{75} 


\god{2002} 


\nomer{11~(486)} %                        Remove in Eng. 


%\nomer{Vol.\,46, No.\,11}%               Set in Eng. 


%\str{\pageref{begin}--\pageref{end}}%  Set in Eng. 


\udk{514.755} 


 


\author{\it Т.Б.\,Жогова} 


% NB:   ^^ remove in Eng. 


\title{Семейства $V^2_{2n-1}$, $U^m_{2n-1}$ и их проективное изгибание} 


 


\date{} 


%\pagestyle{myheadings}%         Set in Eng. 


\pagestyle{plain} %              Remove in Eng. 


\begin{document} 


%\thispagestyle{plain}%          Set in Eng. 


\maketitle 


\label{begin} 


 


 


В данной работе вводятся некоторые классы семейств $L^m_{2n-1}$. 


Семейства $V^2_{2n-1}$ являются обобщением конгруэнций $V$, а семейства 


$U^m_{2n-1}$ --- 


обобщением линейных конгруэнций. Дано конструктивное построение 


семейств $U^m_{2n-1}$. Исследована задача проективного изгибания 


введенных семейств. 


 


Условимся, что по индексам $i$, $j$, $k$ суммирования нет, 


и если они находятся в одном и том же математическом выражении, 


то не принимают равные значения. По индексам $p$, $q$, $r$ всегда 


проводится суммирование. Все индексы, как правило, принимают 


значения от 1 до $n$ включительно. 


\medskip 


 


{\bf 1.} {\it Семейства $V^2_{2n-1}$ и $U^m_{2n-1}$}. В проективном 


пространстве $P_{2n-1}$ 


введем проективный репер $\{A_i,A_{n+i}\}$ с инфинитезимальными 


перемещениями 


$$ 


d A_i=\omega^p_i A_p+\omega^{n+p}_i A_{n+p},\quad 


d A_{n+i}=\omega^p_{n+i}A_p+\omega^{n+p}_{n+i}A_{n+p}, 


$$ 


где формы Пфаффа удовлетворяют известным уравнениям структуры 


проективного пространства. 


 


В [1] введено понятие семейства $L^m_{2n-1}$ и его преобразований Лапласа. 


При этом, если $(n-1)$-плоскость $L_{n-1}=(A_1\dotsc A_n)$ описывает 


семейство $L^m_{2n-1}$, 


то в репере 2-го порядка оно определяется системой уравнений Пфаффа 


\begin{equation} 


\alpha)\ \omega^{n+j}_i=0,\ \ \beta)\ \omega^j_i=a^j_i\omega_j, 


\ \ \gamma)\ \omega^{n+j}_{n+i}=b^j_i\omega_i, 


\end{equation} 


где введено обозначение $\omega^{n+i}_i=\omega_i$, и за независимые формы 


семейства $L^m_{2n-1}$ 


принимаются формы $\omega_k$, $k=\ov{1,m}$, $2\le m\le n$. 


 


В [2] введены абсолютные инварианты $J_{i j}$ и относительные 


инвариантные формы $\varphi_{i j}$ семейства $L^m_{2n-1}$, где 


$$ 


J_{i j}=a^j_i a^i_j/b^j_i b^i_j,\quad 


\varphi_{i j}=\omega^j_i\omega_j+\omega^{n+j}_{n+i}\omega_i. 


$$ 


Нулевые линии формы $\varphi_{i j}$ называются асимптотическими линиями 


фокальной $m$-поверхности $(A_i)$ в направлении фокуса $A_j$. 


 


\begin{defn} 


Семейство $L^2_{2n-1}$ назовем семейством $V^2_{2n-1}$, если 


$$ 


J_{i j}=-1,\quad 


\text{т.\,е.}\quad 


a^j_i a^i_j=-b^j_i b^i_j. 


$$ 





Из определения 1 следует, что $b^j_i=\tau^j_i a^j_i$, $\tau^j_i\tau^i_j=-1$, 


и учитывая $(1\,\gamma))$, получим 


\begin{equation} 


\omega^{n+j}_{n+i}=\tau^j_i a^j_i\omega_i,\quad 


\tau^j_i\tau^i_j=-1. 


\end{equation} 





Замыкая систему уравнений $(1\,\alpha), \beta))$, (2), найдем 


\begin{equation} 


\begin{array}{c} 


a^j_i\Delta a^j_i\wedge \omega_j+\omega^j_{n+i} 


\wedge \omega_i=0,\quad 


\Delta\tau^j_i+\Delta\tau^i_j=0,\\[2mm] 


\tau^j_i a^j_i(\Delta a^j_i+\Delta\tau^j_i 


+A^p_{i j}\omega_p)\wedge \omega_i- 


\omega^j_{n+i}\wedge \omega_j=0,\quad p\ne i,j, 


\end{array} 


\end{equation} 


где 


\begin{gather*} 


\Delta a^j_i=d\ln a^j_i+2\omega^j_j- 


\omega^i_i-\omega^{n+j}_{n+j}- 


(a^j_p a^p_i/a^j_i)\omega_p,\quad p\ne i,j,\\ 


% 


\Delta\tau^j_i=d\ln \tau^j_i+2(\omega^i_i-\omega^{n+i}_{n+i} 


-\omega^j_j+\omega^{n+j}_{n+j}), \\ 


% 


A^k_{i j}=(a^j_k a^k_i/a^j_i) 


(1-\tau^i_j \tau^j_k \tau^k_i). 


\end{gather*} 


\end{defn} 


 


Замкнутая система уравнений $(1\,\alpha), \beta))$, (2), (3), определяющая 


семейство $V^2_{2n-1}$, находится в инволюции с характерами 


$s_1=2n(n-1)$, $s_2=n(n+1)/2-2$, 


т.\,е. справедлива 


 


\begin{thm} 


Семейство $V^2_{2n-1}$ существует с произволом $n(n+1)/2-2$ 


функций двух аргументов. 


\end{thm} 


 


Геометрический смысл уравнения $J_{i j}=-1$ ($i$, $j$ фиксированы) состоит 


в том, что асимптотическим линиям фокальной поверхности $(A_i)$ 


в направлении фокуса $A_j$ соответствует сопряженная система 


линий фокальной поверхности $(A_j)$ в направлении фокуса $A_i$. 


 


\begin{defn} 


Семейство $L^m_{2n-1}$ назовем семейством $U^m_{2n-1}$, если его фокальные 


$m$-по\-верх\-нос\-ти 


вырождаются в прямые. 


\end{defn} 


 


Семейство $U^m_{2n-1}$ при $m\ne n$ определяется замкнутой системой 


уравнений 


\begin{gather} 


\omega^{n+j}_i=0,\ \ \omega^j_i=0, 


\ \ \omega^{n+j}_{n+i}=0,\ \ \omega^j_{n+i}=0, \\ 


% 


\omega_k=\Lambda^r_k\omega_r,\ \ \Delta\Lambda^r_k\wedge\omega_r=0, 


\ \ r=\ov{1,m},\ \ k=\ov{m+1,n}, 


\end{gather} 


где 


$$ 


\Delta\Lambda^j_k=d\Lambda^j_k+\Lambda^j_k 


(\omega^j_j-\omega^{n+j}_{n+j}- 


\omega^k_k+\omega^{n+k}_{n+k}),\quad j=\ov{1,m}. 


$$ 


Эта система находится в инволюции с характерами $s_1=\dotsb=s_m=n-m$, 


т.\,е. справедлива 


 


\begin{thm} 


Семейство $U^m_{2n-1}$ при $m\ne n$ существует с произволом 


$n-m$ функций $m$ аргументов. 


\end{thm} 


 


В силу уравнений (4) 


$$ 


d(A_i A_{n+i})=(\omega^i_i+\omega^{n+i}_{n+i})(A_i A_{n+i}), 


$$ 


т.\,е. прямые $l_i=(A_i A_{n+i})$ неподвижны и, следовательно, 


фокальные $m$-поверхности семейства $L^m_{2n-1}$ вырождаются в прямые 


$l_i$. 


 


Если $m=n$, то уравнения (5) будут отсутствовать, и семейство $U^n_{2n-1}$ 


определится вполне интегрируемой системой (4), т.\,е. справедлива 


 


\begin{thm} 


Семейство $U^n_{2n-1}$ существует с произволом в $4n(n-1)$ постоянных. 


\end{thm} 


 


Система уравнений (4) при фиксированном $i$ является вполне 


интегрируемой, и ее первые интегралы могут быть приняты за 


координаты прямой $l_i$ в пространстве $P_{2n-1}$. Таким образом, чтобы 


конструктивно построить семейство $U^n_{2n-1}$, достаточно в $P_{2n-1}$ 


взять $n$ 


прямых $l_i$ общего положения и рассмотреть множество $(n-1)$-плоскостей, 


каждая из которых пересекает каждую из прямых $l_i$. 


Заметим, что через точку общего положения по отношению к 


прямым $l_i$ проходит единственная $(n-1)$-плоскость семейства 


$U^n_{2n-1}$. 


 


Если $m\ne n$, то нужно еще задать $n-m$ функций $m$ аргументов. Теперь 


формы $\omega_i$ связаны уравнениями (5). Каждое из уравнений 


$\omega_i=0$ является 


вполне интегрируемым, и его первый интеграл является координатой точки 


$A_i$ 


на прямой $l_i$. Зафиксируем в системе (5) индекс $k$. Тогда 


получим уравнение Пфаффа, решение которого существует с 


произволом в одну функцию $m$~аргументов. Пусть $t_i$ --- координата 


точки $A_i$ на прямой $l_i$. Тогда функция 


$t_k=\varphi_k(t_1,\dotsc,t_m)$, определяющая произвол 


решения уравнения (5) при фиксированном $k$, будет индуцировать 


отображение 


\begin{equation} 


q_k:l_1\times l_2\times\dotsb\times l_m\to l_k. 


\end{equation} 


Доказано, что функция $\varphi_k$ однозначно разрешима относительно 


любого из аргументов $t_1,\dotsc,t_m$. В случае $m=2$, $n=3$ отображение 


(6) определяет квазигруппу, которая была исследована в [3]. 


\medskip 


 


{\bf 2.} {\it Изгибание семейств $V^2_{2n-1}$ и $U^m_{2n-1}$}. 


Рассматривается проективное изгибание Фубини--Картана ([4], с.\,102). 


Отнесем семейство $\ov L{}^m_{2n-1}$, на которое проективным изгибанием 


2-го 


порядка наложимо семейство $L^m_{2n-1}$, к реперу $\{B_i,B_{n+i}\}$ с 


инфинитезимальными перемещениями 


$$ 


d B_i=\Omega^q_i B_q+\Omega^{n+q}_i B_{n+q},\quad 


d B_{n+i}=\Omega^q_{n+i}B_q+\Omega^{n+q}_{n+i}B_{n+q}. 


$$ 


Если $(n-1)$-плоскость $\ov L_{n-1}=(B_1\dotsc B_n)$ описывает семейство 


$\ov L{}^m_{2n-1}$, а $\Pi$ --- такое 


проективное преобразование, что 


$$ 


\Pi(B_i)=\wt B_i=A_i,\quad 


\Pi(B_{n+i})=\wt B_{n+i}=A_{n+i},\quad 


\Pi(\ov L{}^m_{2n-1})=\wt L{}^m_{2n-1}, 


$$ 


то условие касания 2-го порядка семейств $L^m_{2n-1}$ и $\wt 


L{}^m_{2n-1}$ определится уравнением 


$$ 


\wt L_{n-1}+d\wt L_{n-1}+\tfrac{1}{2}d^2\wt L_{n-1}= 


(\theta_0+\theta_1+\tfrac{1}{2}\theta_2)(L_{n-1}+d L_{n-1}+ 


\tfrac{1}{2}d^2L_{n-1}). 


$$ 


 


В дальнейшем используются обозначения 


$$ 


\Omega^{n+i}_i=\Omega_i,\quad 


\wt\omega{}^\beta_\alpha=\Omega^\beta_\alpha-\omega^\beta_\alpha, 


\quad \alpha,\beta=\ov{1,2n}. 


$$ 


 


Проективное изгибание 1-го порядка семейства $V^2_{2n-1}$ определяется 


системой $(1\,\alpha), \beta))$, (2), (3), к которой добавляются уравнения 


\begin{gather} 


\wt\omega_i=0,\quad \Omega^{n+j}_i=0,\\ 


%


(\wt\omega{}^i_i-\wt\omega{}^{n+i}_{n+i})\wedge\omega_i=0,\quad 


\wt\omega{}^j_i\wedge\omega_j-\wt\omega{}^{n+j}_{n+i} 


\wedge\omega_i=0. 


\end{gather} 


Исследование замкнутой системы уравнений $(1\,\alpha), \beta))$, (2), (3), 


(7), (8) показывает, что справедлива 


 


\begin{thm} 


Существует с произволом в $n(n-1)$ функций двух аргументов 


семейство $\ov L{}^2_{2n-1}$, на которое изгибанием $1$-го порядка с 


произволом в $n$ функций одного аргумента наложимо заданное семейство 


$V^2_{2n-1}$. 


\end{thm} 


 


Изгибание 1-го порядка семейства $U^m_{2n-1}$ определяется замкнутой 


системой уравнений (4), (5), (7), (8), исследование которой приводит


к следующей теореме.


 


\begin{thm} 


Существует с произволом в $n(n-1)$ функций двух аргументов 


семейство $\ov L{}^m_{2n-1}$, на которое изгибанием $1$-го порядка с 


произволом в $n$ функций одного аргумента наложимо заданное семейство 


$U^m_{2n-1}$. 


\end{thm} 


 


\begin{defn} 


Фокальным изгибанием $k$-го порядка семейства $L^m_{2n-1}$ 


назовем изгибание $k$-го порядка его фокальных $m$-поверхностей. 


\end{defn} 


 


При фокальном изгибании 1-го порядка семейства $L^m_{2n-1}$ имеют


место уравнения 


$$ 


\wt B_i+d\wt B_i=(\theta^i_0+\theta^i_1) 


(A_i+d A_i). 


$$ 


 


Рассматривая фокальное изгибание 1-го порядка семейства $V^2_{2n-1}$, 


получим уравнения (7) и уравнения 


\begin{equation} 


\wt\omega{}^j_i=0, 


\end{equation} 


т.\,е. фокальное изгибание 1-го порядка семейства $V^2_{2n-1}$ влечет 


за собой его изгибание 1-го порядка. Система (9) с учетом $(1\,\beta))$ 


имеет дифференциальные следствия 


\begin{equation} 


\wt\omega{}^i_i+\wt\omega{}^j_j= 


p_i\omega_i-p_j\omega_j, 


\end{equation} 


а замыкание системы (7), (9), (10) имеет вид 


\begin{gather} 


(\wt\omega{}^i_i-\wt\omega{}^{n+i}_{n+i}) 


\wedge\omega_i=0,\quad \wt\omega{}^{n+j}_{n+i} 


\wedge\omega_i=0, \\ 


%


(\wt\omega{}^j_{n+i}+p_i a^j_i\omega_j) 


\wedge\omega_i=0,\quad \Delta p_i\wedge\omega_i=0, 


\end{gather} 


где 


$$ 


\Delta p_i=d p_i+p_i(\omega^i_i-\omega^{n+i}_{n+i}) 


+\wt\omega{}^i_{n+i}. 


$$ 


Исследование замкнутой системы уравнений $(1\,\alpha), \beta))$, (2), (3),


(7), (9)--(12) показывает, что справедлива 


 


\begin{thm} 


Существует с произволом в $n(2n-1)$ функций одного аргумента 


семейство $\ov L{}^2_{2n-1}$, на которое фокальным изгибанием $1$-го 


порядка 


с произволом в $n$ функций одного аргумента наложимо заданное 


семейство $V^2_{2n-1}$. 


\end{thm} 


 


При фокальном изгибании 1-го порядка семейства $U^m_{2n-1}$ 


имеют место уравнения (7), (9), внешним дифференцированием 


которых получим уравнения (11) и уравнения 


\begin{equation} 


\wt\omega{}^j_{n+i}\wedge\omega_i=0. 


\end{equation} 


 


Исследование замкнутой системы уравнений (4), (5), (7), (9), 


(11), (13) показывает, что имеет место 


 


\begin{thm} 


Существует с произволом в $2n(n-1)$ функций одного аргумента 


семейство $\ov L{}^m_{2n-1}$, на которое фокальным изгибанием $1$-го 


порядка с произволом в $n$ функций одного аргумента 


наложимо заданное семейство $U^m_{2n-1}$. 


\end{thm} 


 


Рассмотрим изгибание 1-го порядка семейства $V^2_{2n-1}$ одновременно 


с изгибанием 1-го порядка его прямых Лапласа. Тогда дополнительно 


будут иметь место условия 


$$ 


(\wt B_i\wt B_{n+i})+d(\wt B_i\wt B_{n+i})= 


(\wt\theta{}^i_0+\wt\theta{}^i_1)((A_i A_{n+i})+ 


d(A_i A_{n+i})). 


$$ 


Оказалось, что в этом случае семейство $V^2_{2n-1}$ является семейством 


$U^2_{2n-1}$, а семейство $\ov L{}^2_{2n-1}$ является семейством 


$\ov U{}^2_{2n-1}$ и справедлива 


 


\begin{thm} 


Заданное семейство $U^2_{2n-1}$ допускает изгибание $1$-го порядка 


одновременно с прямыми Лапласа с произволом в $n$ 


функций одного аргумента. 


\end{thm} 


 


При изгибании 2-го порядка семейства $U^m_{2n-1}$ справедлива 


 


\begin{thm} 


Заданное семейство $U^m_{2n-1}$ допускает изгибание $2$-го порядка 


с произволом в~$n$ функций одного аргумента. 


\end{thm} 


 


\begin{thebibliography}{9} 


 


\bibitem{1} 


Макеев Г.Н. {\em О некотором обобщении преобразований Лапласа} 


// Изв. вузов. Математика. -- 1975. -- \No\,2. -- С.\,123--125. 


\bibitem{2} 


Макеев Г.Н. {\em Семейства $R^m_{2n-1}$ и $\Phi^n_{2n-1}$} // Изв. вузов. 


Математика. -- 1990. -- \No\,1. -- С.\,84--86. 


\bibitem{3} 


Жогова Т.Б. {\em О квазигруппе, порождаемой одним классом 


двупараметрических семейств двумерных плоскостей в $P_5$} // 


Изв. вузов. Математика. -- 1980. -- \No\,2. -- С.\,63--66. 


\bibitem{4} 


Фиников С.П. {\em Теория пар конгруэнций}. -- М.: ГИТТЛ, 1956. -- 443 с. 


 


\end{thebibliography} 


 


\vskip 2cm 


 


\post{\it Нижегородский государственный}{\it Поступила} 


\post{\it педагогический университет}{17.07.2001} 


 


\label{end} 


\end{document} 


